\section{习题}

\subsection{本章速览}

\begin{itemize}
\item \textbf{费马引理：}若 $f$ 在 $x_0$ 处取得极值且在该点可导，则 $f'(x_0)=0$。反之不成立：$f(x)=x^3$ 在 $x=0$ 处有 $f'(0)=0$，但 $x=0$ 不是极值点。
\item \textbf{罗尔定理：}三条件缺一不可——$f$ 在 $[a,b]$ 上连续、在 $(a,b)$ 内可导、$f(a)=f(b)$；结论是存在 $\xi\in(a,b)$ 使 $f'(\xi)=0$。只给开区间内可导而不给端点连续，结论可能失效。
\item \textbf{拉格朗日中值定理：}$f(b)-f(a)=f'(\xi)(b-a)$，$\xi\in(a,b)$。它把函数增量、自变量增量与导数联系起来，是证明不等式、判断单调性和推出“导数恒为零则函数为常数”的基本工具。
\item \textbf{柯西中值定理：}在连续、可导之外还要求 $g'(x)\neq 0$，结论为 $\frac{f(b)-f(a)}{g(b)-g(a)}=\frac{f'(\xi)}{g'(\xi)}$。取 $g(x)=x$ 即得拉格朗日中值定理。
\item \textbf{泰勒公式：}带皮亚诺余项 $o((x-x_0)^n)$ 用于求极限、判断无穷小阶数；带拉格朗日余项 $\dfrac{f^{(n+1)}(\xi)}{(n+1)!}(x-x_0)^{n+1}$ 用于不等式证明与误差估计。两种余项的多项式部分完全相同。
\item \textbf{辅助函数套路：}要证含 $f'(\xi)$ 与 $f(\xi)$ 的关系式，把结论中的 $\xi$ 换成 $x$ 看成微分方程，其解就是辅助函数。常用：$e^{g(x)}f(x)$、$x^nf(x)$、$\dfrac{f(x)}{g(x)}$、$f(x)g(x)$。
\end{itemize}

\subsection{基础过关题}

\textbf{1.} 验证函数 $f(x)=x^{3}-9x^{2}+23x-15$ 在区间 $[1,5]$ 上满足罗尔定理的条件，并求出定理结论中所有的 $\xi$。

\ifshowsolutions%
\boxed{\textbf{解}}
\textbf{思路：}先验证三个条件，再解方程 $f'(x)=0$，最后检查所得根是否落在开区间 $(1,5)$ 内。

$f$ 是多项式，在 $[1,5]$ 上连续、在 $(1,5)$ 内可导。计算端点值：
\[
f(1)=1-9+23-15=0,\qquad f(5)=125-225+115-15=0,
\]
故 $f(1)=f(5)$，罗尔定理的三个条件全部满足。

求导并因式分解：
\[
f'(x)=3x^{2}-18x+23.
\]
令 $f'(x)=0$，得
\[
x=\frac{18\pm\sqrt{324-4\cdot 3\cdot 23}}{6}
=\frac{18\pm\sqrt{48}}{6}
=3\pm\frac{2\sqrt3}{3}.
\]
其中 $3-\frac{2\sqrt3}{3}\approx 1.845$，$3+\frac{2\sqrt3}{3}\approx 4.155$，二者都落在 $(1,5)$ 内，因此都是罗尔定理中的 $\xi$。

\boxed{\textbf{答案：}\xi=3\pm\dfrac{2\sqrt3}{3}}
\else \answerarea[8] \fi%

\textbf{2.} 设 $f(x)=\ln x$。验证 $f$ 在 $[1,e]$ 上满足拉格朗日中值定理的条件，并求出定理中的 $\xi$。

\ifshowsolutions%
\boxed{\textbf{解}}
\textbf{思路：}$\ln x$ 在 $[1,e]$ 上连续、在 $(1,e)$ 内可导，直接由中值公式列方程解 $\xi$。

$f(x)=\ln x$ 在 $[1,e]$ 上连续，在 $(1,e)$ 内可导，且 $f'(x)=\dfrac1x$，条件满足。由拉格朗日中值定理，存在 $\xi\in(1,e)$ 使
\[
\frac{f(e)-f(1)}{e-1}=f'(\xi),
\]
即
\[
\frac{1-0}{e-1}=\frac1{\xi}.
\]
解得 $\xi=e-1$。由于 $e\approx 2.718$，故 $\xi=e-1\approx 1.718$，确实满足 $1<\xi<e$。

\boxed{\textbf{答案：}\xi=e-1}
\else \answerarea[8] \fi%

\textbf{3.} 设 $f(x)=\sin x$，$g(x)=\cos x$。在区间 $\left[0,\dfrac{\pi}{2}\right]$ 上验证柯西中值定理，并求出定理中的 $\xi$。

\ifshowsolutions%
\boxed{\textbf{解}}
\textbf{思路：}柯西中值定理比对值 $\dfrac{f'(\xi)}{g'(\xi)}$，先确认 $g'$ 在开区间内不为零，再解三角方程。

$f,g$ 在 $\left[0,\frac{\pi}{2}\right]$ 上连续，在 $\left(0,\frac{\pi}{2}\right)$ 内可导，且
\[
g'(x)=-\sin x\neq 0,\qquad x\in\left(0,\frac{\pi}{2}\right),
\]
条件满足。计算两端之差：
\[
f\left(\frac{\pi}{2}\right)-f(0)=1-0=1,\qquad
g\left(\frac{\pi}{2}\right)-g(0)=0-1=-1.
\]
由柯西中值定理，存在 $\xi\in\left(0,\frac{\pi}{2}\right)$ 使
\[
\frac{1}{-1}=\frac{\cos\xi}{-\sin\xi}=-\cot\xi,
\]
即 $\cot\xi=1$，从而 $\tan\xi=1$，$\xi=\dfrac{\pi}{4}$，它落在 $\left(0,\frac{\pi}{2}\right)$ 内。

\boxed{\textbf{答案：}\xi=\dfrac{\pi}{4}}
\else \answerarea[8] \fi%

\textbf{4.} 证明：当 $x>0$ 时，
\[
\frac{x}{1+x}<\ln(1+x)<x.
\]

\ifshowsolutions%
\boxed{\textbf{解}}
\textbf{思路：}把 $\ln(1+x)-\ln 1$ 写成 $x\cdot\dfrac{1}{1+\xi}$，再用 $\xi\in(0,x)$ 夹出上下界。

对 $F(t)=\ln(1+t)$ 在 $[0,x]$ 上使用拉格朗日中值定理。$F$ 在 $[0,x]$ 上连续、在 $(0,x)$ 内可导，$F'(t)=\dfrac{1}{1+t}$，故存在 $\xi\in(0,x)$ 使
\[
\ln(1+x)-\ln 1=\frac{1}{1+\xi}\,(x-0),
\]
即
\[
\ln(1+x)=\frac{x}{1+\xi}.
\]
由 $0<\xi<x$ 得 $1<1+\xi<1+x$，取倒数（各项为正）得
\[
\frac{1}{1+x}<\frac{1}{1+\xi}<1.
\]
三边同乘正数 $x$，得
\[
\frac{x}{1+x}<\frac{x}{1+\xi}<x,
\]
即 $\dfrac{x}{1+x}<\ln(1+x)<x$。

\boxed{\textbf{答案：}\text{不等式对一切 }x>0\text{ 成立}}
\else \answerarea[8] \fi%

\textbf{5.} 求极限
\[
\lim_{x\to 0}\frac{\cos x-e^{-x^{2}/2}}{x^{4}}.
\]

\ifshowsolutions%
\boxed{\textbf{解}}
\textbf{思路：}分子中两项都展开到 $x^4$ 阶，低阶项应恰好抵消，剩余的最低阶项决定极限。

由麦克劳林公式
\[
\cos x=1-\frac{x^{2}}{2}+\frac{x^{4}}{24}+o(x^{4}).
\]
对 $e^{u}=1+u+\dfrac{u^{2}}{2}+o(u^{2})$ 取 $u=-\dfrac{x^{2}}{2}$，得
\[
e^{-x^{2}/2}=1-\frac{x^{2}}{2}+\frac{1}{2}\left(-\frac{x^{2}}{2}\right)^{2}+o(x^{4})
=1-\frac{x^{2}}{2}+\frac{x^{4}}{8}+o(x^{4}).
\]
两式相减，$1$ 与 $-\dfrac{x^{2}}{2}$ 都抵消：
\[
\cos x-e^{-x^{2}/2}
=\left(\frac{1}{24}-\frac18\right)x^{4}+o(x^{4})
=-\frac{1}{12}x^{4}+o(x^{4}).
\]
于是
\[
\lim_{x\to 0}\frac{\cos x-e^{-x^{2}/2}}{x^{4}}
=\lim_{x\to 0}\frac{-\frac{1}{12}x^{4}+o(x^{4})}{x^{4}}
=-\frac{1}{12}.
\]

\boxed{\textbf{答案：}-\dfrac{1}{12}}
\else \answerarea[8] \fi%

\subsection{综合提高题}

\textbf{6.} 讨论方程 $x^{3}-3x+1=0$ 的实根个数，并指出每个根所在的区间。

\ifshowsolutions%
\boxed{\textbf{解}}
\textbf{思路：}先用导数确定单调区间与极值，结合两端极限说明“至少”有几个根；再用罗尔定理说明“至多”有几个根。

设 $f(x)=x^{3}-3x+1$，则
\[
f'(x)=3x^{2}-3=3(x-1)(x+1).
\]
$f'$ 的零点为 $x=-1$ 和 $x=1$，列表可知：$f$ 在 $(-\infty,-1)$ 上递增，在 $(-1,1)$ 上递减，在 $(1,+\infty)$ 上递增。极值为
\[
f(-1)=-1+3+1=3>0\ (\text{极大值}),\quad
f(1)=1-3+1=-1<0\ (\text{极小值}).
\]
又 $\lim\limits_{x\to-\infty}f(x)=-\infty$，$\lim\limits_{x\to+\infty}f(x)=+\infty$。

在三个单调区间上分别用零点定理：
\begin{itemize}
\item $(-\infty,-1)$：$f(-2)=-8+6+1=-1<0$，$f(-1)=3>0$，故在 $(-2,-1)$ 内有一个根，且由严格单调知只有一个；
\item $(-1,1)$：$f(-1)=3>0$，$f(1)=-1<0$，故在 $(-1,1)$ 内有一个根，且只有一个；
\item $(1,+\infty)$：$f(1)=-1<0$，$f(2)=8-6+1=3>0$，故在 $(1,2)$ 内有一个根，且只有一个。
\end{itemize}
所以至少有 $3$ 个实根。若实根多于 $3$ 个，即至少有 $4$ 个，则对 $f$ 连续使用三次罗尔定理，$f'$ 至少要有 $3$ 个零点；但 $f'=3(x-1)(x+1)$ 只有 $x=\pm1$ 两个零点，矛盾。因此恰有 $3$ 个实根。

\boxed{\textbf{答案：}\text{恰有 }3\text{ 个实根，分别在 }(-2,-1),\,(-1,1),\,(1,2)\text{ 内}}
\else \answerarea[10] \fi%

\textbf{7.} 设 $f$ 在 $[0,1]$ 上连续，在 $(0,1)$ 内可导，$f(0)=0$，且 $f'$ 在 $(0,1)$ 内严格单调递增。证明：函数
\[
g(x)=\frac{f(x)}{x}
\]
在 $(0,1)$ 内严格单调递增。

\ifshowsolutions%
\boxed{\textbf{解}}
\textbf{思路：}直接对 $g$ 求导，分子 $xf'(x)-f(x)$ 可以借助中值定理写成 $x\bigl(f'(x)-f'(\xi)\bigr)$，再用 $f'$ 的单调性定号。

对 $x\in(0,1)$，
\[
g'(x)=\frac{xf'(x)-f(x)}{x^{2}}.
\]
因为 $f$ 在 $[0,x]$ 上连续、在 $(0,x)$ 内可导，由拉格朗日中值定理，存在 $\xi\in(0,x)$ 使
\[
f(x)-f(0)=f'(\xi)(x-0),
\]
结合 $f(0)=0$ 得 $f(x)=xf'(\xi)$。代入分子：
\[
xf'(x)-f(x)=xf'(x)-xf'(\xi)=x\bigl[f'(x)-f'(\xi)\bigr].
\]
由于 $0<\xi<x<1$ 且 $f'$ 在 $(0,1)$ 内严格单调递增，故 $f'(x)>f'(\xi)$，从而
\[
xf'(x)-f(x)>0.
\]
于是 $g'(x)>0$ 对一切 $x\in(0,1)$ 成立，$g$ 在 $(0,1)$ 内严格单调递增。

\boxed{\textbf{答案：}\text{证明如上。}}
\else \answerarea[10] \fi%

\textbf{8.} 证明：当 $x>0$ 时，
\[
\ln(1+x)<x-\frac{x^{2}}{2}+\frac{x^{3}}{3}.
\]

\ifshowsolutions%
\boxed{\textbf{解}}
\textbf{思路：}右端正是 $\ln(1+x)$ 的三阶泰勒多项式。把函数在 $0$ 处展开到三阶，用拉格朗日余项的符号判断大小。

设 $f(t)=\ln(1+t)$，则
\[
f(0)=0,\quad f'(t)=\frac{1}{1+t},\quad f''(t)=-\frac{1}{(1+t)^{2}},
\]
\[
f'''(t)=\frac{2}{(1+t)^{3}},\quad f^{(4)}(t)=-\frac{6}{(1+t)^{4}}.
\]
于是 $f'(0)=1$，$f''(0)=-1$，$f'''(0)=2$。在 $[0,x]$ 上使用带拉格朗日余项的三阶泰勒公式，存在 $\xi\in(0,x)$ 使
\[
f(x)=f(0)+f'(0)x+\frac{f''(0)}{2!}x^{2}+\frac{f'''(0)}{3!}x^{3}
+\frac{f^{(4)}(\xi)}{4!}x^{4}.
\]
代入各导数：
\[
\ln(1+x)=x-\frac{x^{2}}{2}+\frac{x^{3}}{3}
-\frac{6}{(1+\xi)^{4}}\cdot\frac{x^{4}}{24}
=x-\frac{x^{2}}{2}+\frac{x^{3}}{3}-\frac{x^{4}}{4(1+\xi)^{4}}.
\]
因为 $x>0$，$\xi\in(0,x)$，所以 $\dfrac{x^{4}}{4(1+\xi)^{4}}>0$，从而
\[
\ln(1+x)<x-\frac{x^{2}}{2}+\frac{x^{3}}{3}.
\]

\boxed{\textbf{答案：}\text{证明如上。}}
\else \answerarea[10] \fi%

\textbf{9.} 设 $f$ 在 $[0,1]$ 上具有连续的二阶导数，$f(0)=f(1)=0$，且 $|f''(x)|\le M$（$x\in[0,1]$）。证明：对任意 $x\in[0,1]$，
\[
|f'(x)|\le\frac{M}{2}.
\]

\ifshowsolutions%
\boxed{\textbf{解}}
\textbf{思路：}在任意点 $x$ 处向两个端点各写一次带拉格朗日余项的一阶泰勒公式，两式相减消去 $f(x)$，只剩 $f'(x)$ 与两个二阶导数。

先设 $x\in(0,1)$。分别以 $x$ 为展开点，把 $f(0)$、$f(1)$ 展开：
\[
0=f(0)=f(x)+f'(x)(0-x)+\frac{f''(\xi_1)}{2}(0-x)^{2},
\]
\[
0=f(1)=f(x)+f'(x)(1-x)+\frac{f''(\xi_2)}{2}(1-x)^{2},
\]
其中 $\xi_1\in(0,x)$，$\xi_2\in(x,1)$。即
\[
0=f(x)-xf'(x)+\frac{x^{2}}{2}f''(\xi_1),\tag{1}
\]
\[
0=f(x)+(1-x)f'(x)+\frac{(1-x)^{2}}{2}f''(\xi_2).\tag{2}
\]
$(1)-(2)$ 消去 $f(x)$：
\[
0=-f'(x)+\frac{x^{2}}{2}f''(\xi_1)-\frac{(1-x)^{2}}{2}f''(\xi_2),
\]
故
\[
f'(x)=\frac{x^{2}}{2}f''(\xi_1)-\frac{(1-x)^{2}}{2}f''(\xi_2).
\]
取绝对值并用 $|f''|\le M$：
\[
|f'(x)|\le\frac{M}{2}\bigl[x^{2}+(1-x)^{2}\bigr].
\]
而
\[
x^{2}+(1-x)^{2}=2\left(x-\frac12\right)^{2}+\frac12\le 1,\qquad x\in[0,1],
\]
所以 $|f'(x)|\le\dfrac{M}{2}$。

再补端点。对 $x=0$，在 $[0,1]$ 上把 $f(1)$ 在 $0$ 处展开：
\[
0=f(1)=f(0)+f'(0)\cdot 1+\frac{f''(\xi)}{2}\cdot 1^{2}
=f'(0)+\frac{f''(\xi)}{2},
\]
故 $|f'(0)|=\dfrac{|f''(\xi)|}{2}\le\dfrac{M}{2}$。同理，把 $f(0)$ 在 $1$ 处展开可得 $|f'(1)|\le\dfrac{M}{2}$。

最后说明系数 $\dfrac12$ 不可改进：取 $f(x)=\dfrac{M}{2}x(1-x)$，则 $f(0)=f(1)=0$，$f''(x)=-M$ 满足 $|f''|\le M$，而 $f'(0)=\dfrac{M}{2}$，恰好取到等号。

\boxed{\textbf{答案：}|f'(x)|\le\dfrac{M}{2}\text{，且 }\dfrac12\text{ 不可改进}}
\else \answerarea[10] \fi%

\subsection{真题风格与综合训练}

\textbf{10.}（单项选择题）下列命题中正确的是

\begin{center}
\begin{tabular}{p{0.46\linewidth}p{0.46\linewidth}}
(A) 若 $f$ 在 $x_0$ 处取得极值，则 $f'(x_0)=0$； & (B) 若 $f$ 在 $[a,b]$ 上连续、在 $(a,b)$ 内可导且 $f(a)=f(b)$，则存在 $\xi\in(a,b)$ 使 $f'(\xi)=0$；\\[4pt]
(C) 若 $f'(\xi)=0$，则 $\xi$ 必为 $f$ 的极值点； & (D) 若 $f$ 在 $[a,b]$ 上连续、在 $(a,b)$ 内可导且只有一个驻点，则该驻点必为极值点。
\end{tabular}
\end{center}

\ifshowsolutions%
\boxed{\textbf{解}}
\textbf{思路：}逐项检查定理条件，凡缺少“可导”“连续”或把必要条件当充分条件者均可构造反例。

(A) 错。费马引理还要求 $f$ 在 $x_0$ 处可导。反例：$f(x)=|x|$ 在 $x=0$ 处取得极小值，但 $f'(0)$ 不存在。

(B) 对。这正是罗尔定理，三个条件齐备。

(C) 错。驻点未必是极值点。反例：$f(x)=x^{3}$，$f'(0)=0$，但 $f$ 在 $0$ 两侧都递增，$x=0$ 不是极值点。

(D) 错。反例：$f(x)=x^{3}$ 在 $[-1,1]$ 上只有一个驻点 $x=0$，但 $x=0$ 不是极值点。

\boxed{\textbf{答案：}(B)}
\else \answerarea[8] \fi%

\textbf{11.}（填空题）当 $x\to 0$ 时，$\ln(1+x)-x+\dfrac{x^{2}}{2}$ 与 $ax^{k}$ 是等价无穷小，则 $a=\underline{\qquad}$，$k=\underline{\qquad}$。

\ifshowsolutions%
\boxed{\textbf{解}}
\textbf{思路：}把 $\ln(1+x)$ 展开到首个未被抵消的阶，比较系数即可。

由麦克劳林公式
\[
\ln(1+x)=x-\frac{x^{2}}{2}+\frac{x^{3}}{3}+o(x^{3}),
\]
所以
\[
\ln(1+x)-x+\frac{x^{2}}{2}=\frac{x^{3}}{3}+o(x^{3}).
\]
因此它与 $\dfrac13x^{3}$ 等价，即 $a=\dfrac13$，$k=3$。

\boxed{\textbf{答案：}a=\dfrac13,\ k=3}
\else \answerarea[3] \fi%

\textbf{12.}（解答题）设 $f$ 在 $[0,1]$ 上具有连续的二阶导数，$f(0)=f(1)=0$，$f'(0)=1$。证明：存在 $\xi\in(0,1)$，使得
\[
f''(\xi)=-2.
\]

\ifshowsolutions%
\boxed{\textbf{解}}
\textbf{思路：}要证的式子等价于某个辅助函数的二阶导数为零。把 $-2$ 移过去看成 $f''+2=0$，积分两次得辅助函数 $g(x)=f(x)+x^{2}-x$。

令
\[
g(x)=f(x)+x^{2}-x=f(x)-x(1-x).
\]
则 $g$ 在 $[0,1]$ 上具有连续的二阶导数，且
\[
g(0)=f(0)+0-0=0,\qquad g(1)=f(1)+1-1=0.
\]
由罗尔定理，存在 $c\in(0,1)$ 使
\[
g'(c)=0.
\]
另一方面
\[
g'(x)=f'(x)+2x-1,\qquad g'(0)=f'(0)-1=1-1=0.
\]
于是 $g'$ 在 $[0,c]$ 上连续、在 $(0,c)$ 内可导，且 $g'(0)=g'(c)=0$。再次使用罗尔定理，存在 $\xi\in(0,c)\subset(0,1)$ 使
\[
g''(\xi)=0.
\]
而 $g''(x)=f''(x)+2$，故 $f''(\xi)+2=0$，即 $f''(\xi)=-2$。

\boxed{\textbf{答案：}\text{证明如上。}}
\else \answerarea[10] \fi%
